\section{Aligned splits for the residual witnesses}\label{sec:aligned}

The two outer splits leave the task of computing $F_2$ on $D_0^*$.
Every episode in $D_0$ has one of the five witnesses from
Lemma~\ref{lem:witness-coverage}. We eliminate them in order. At stage $i$,
the current computation is a map $g_{i,e}:U_i\to U_i$, with
$g_{0,e}=F_{2,e}$ and $U_0=V\oplus\mathbf F_2$. We construct a new affine
update $g_{i+1,e}$ on a larger vector space. Two splits compute this new
update while removing the current witness; its linear part then recovers
$g_{i,e}$, by a basis encoding described below. The new update is the
same for both splits and is defined on every complete episode.

The witnesses serve to decode the cut position. A suffix tests every
candidate origin, excludes later witnesses there, and requires a unique
candidate with the current witness. After decoding, the first split uses
a perfectly predicted trial. On the remaining domain, a second split
uses a full decision tree to recover a prefix product from one of its
measurements. The secondary guard must exclude later witnesses on a
larger set of translates: this will ensure that the actual history has
no perfectly predicted trial. Proposition~\ref{prop:parameters} supplies
the finite schedules and ambiguity bounds used here.

\subsection{Stencils, domains, and the induction invariant}

For $i=0,\ldots,4$, place a nonempty finite stencil $\mathcal P_i$ of cut
offsets from $s$ in its own subbuffer. The buffers occur in the reverse
order $4,3,2,1,0$, after $s$ and before $t$, and the $i$-th ends at $a_i$.
Thus $a_{i+1}$ is before the entire stage-$i$ buffer whenever $i<4$.

\begin{figure}[htbp]
\centering
\begin{tikzpicture}[x=0.83cm,y=0.85cm,>=Stealth,font=\small]
 \draw[->] (0,0)--(13,0);
 \foreach \x/\label in {0/s,7/t,11/q,12.5/{q+B_{\rm tail}}}
   {\draw (\x,-.10)--(\x,.10);
    \node[below=4pt] at (\x,0) {$\label$};}
 \foreach \x/\i in {0.4/4,1.65/3,2.9/2,4.15/1,5.4/0}
   {\draw[fill=black!5] (\x,.4) rectangle ++(1,.55);
    \node at (\x+.5,.675) {$i=\i$};
    \draw (\x+1,.1)--(\x+1,.35);
    \node[above,font=\scriptsize] at (\x+1,1.02) {$a_{\i}$};}
 \draw[decorate,decoration={brace,amplitude=4pt}]
   (7.3,1.05)--(10.3,1.05);
 \node[above=7pt,align=center] at (8.8,1.05) {outer trials\\and windows};
 \draw[fill=black!5] (7.3,.4) rectangle (10.3,.95);
 \draw[dashed] (10.6,.4)--(11,.4);
 \node[above,font=\scriptsize,align=center] at (11.4,1.0)
       {end rule\\and padding};
\end{tikzpicture}
\caption{The aligned buffers occur in reverse stage order. For $i<4$, the
stage-$i$ buffer and the entire suffix beginning at $a_i$ lie after
$a_{i+1}$. Their data are therefore on the reading side of a suffix
starting at $a_{i+1}$. The drawing is schematic; the end-marker lag
and padding are chosen after every bounded inspection range.}\label{fig:reverse-buffers}
\end{figure}

Write
\begin{equation}\label{eq:gamma}
 \Delta_i=\mathcal P_i-\mathcal P_i,\qquad
 \Gamma_i=\sum_{j<i}(\Delta_j+\Delta_j),\qquad
 \Gamma_0=\{0\}.
\end{equation}
Every $\Delta_i,\Gamma_i$ contains zero, and
$\Gamma_{i+1}=\Gamma_i+\Delta_i+\Delta_i$.
For any finite set $\Gamma$, define
\[
 A_h^\Gamma(t,q)
 \quad\Longleftrightarrow\quad
 \text{there is }\gamma\in\Gamma\text{ with }A_h(t+\gamma,q).
\]
The anchor $q$ remains fixed. The entry domain $D_i\subseteq D_0$ has
the promise
\begin{equation}\label{eq:entry-promise}
                 \bigvee_{h\ge i}A_h^{\Gamma_i}(t,q).
\end{equation}
For $i=0$, this is Lemma~\ref{lem:witness-coverage}; the two full sets
defined below will give the induction step. Graph tests at stage $i$
concern $g_i$ on $D_i^*$, although every $g_{i,e}$ is defined on all
complete episodes.

The quantitative hypothesis is
\begin{equation}\label{eq:ambiguity-bound}
 \#\{(p,p')\in\mathcal P_i^2:p\ne p',\
                A_i^{\Gamma_i}(t+p-p',q)\}\le C_i.
\end{equation}
Every required kind of cut has more than $C_i$ copies in $\mathcal P_i$.
Also $h_i=\diam\mathcal P_i<w_E$ for $i<4$.
All witness controls at every candidate and every required translate
are after the cut and inside the true first episode. A suffix must
itself contain their full reading intervals to pass.

The reverse order of the buffers has a computational purpose. From
$a_{i+1}$, the whole stage-$i$ buffer and the suffix starting at $a_i$
are still unread. This lets a suffix formula for $g_i$ be incorporated
into a suffix formula for $g_{i+1}$, as shown next.



\subsection{Total tables and the new affine update}

The suffix beginning at $a_i$ does not know the product $T_{s,a_i}$ of
the omitted prefix. We therefore keep a table for every possible value
of that product. For fixed suffix letters and fixed $t,q$, the table is
a finite formula
\[
                 \Phi_i(m):U_i\longrightarrow U_i\qquad(m\in M)
\]
with
\begin{equation}\label{eq:table-property}
                      \Phi_i(T_{s,a_i})=g_{i,e}.
\end{equation}
It is defined for every $m$, including values that no prefix realizes.
It may use the indicated suffix, bounded marker data, $M$-products,
and fixed $t,q$. Evaluating this table is a finite calculation on one
episode; it is not a query to a graph language for a sequence of episodes.
We prove the existence of these tables inductively below.

To carry this table through another pair of splits, encode a payload
$v\in U_i$ and a work-register value $m\in M$ as a single basis vector.
Set
\[
                   U_{i+1}=\mathbf F_2^{\,U_i\times M}
\]
with basis $[v,m]$. A letter product $d\in M$ acts linearly by
$[v,m]\rho(d)=[v,md]$. At the end of the stage-$i$ buffer, replace the
remaining computation by the terminal linear map $H_i$ defined by
\begin{equation}\label{eq:late-map}
                    [v,m]H_i=[\Phi_i(m)(v),1_M].
\end{equation}
This is a linear extension from the displayed basis assignment; it
does not require $\Phi_i(m)$ to be linear on $U_i$. The unshifted
episode computation consists of $\rho(T_{s,a_i})$ followed by $H_i$.
No letter action follows $H_i$, which already represents the entire
remaining suffix. In particular,
\[
 [v,1_M]\rho(T_{s,a_i})H_i
   =[\Phi_i(T_{s,a_i})(v),1_M]
   =[g_{i,e}(v),1_M].
\]
Resetting the work register to $1_M$ makes the output suitable as the
input to the next episode. If $L_e=\rho(T_{s,a_i})H_i$, then for two
episodes $e,f$,
\[
                   [v,1_M]L_eL_f
                      =[g_{i,f}(g_{i,e}(v)),1_M].
\]
Thus products of these linear maps recover the composed $g_i$ update.
Our new affine update will have precisely $L_e$ as its linear part.

In the first part of the stage-$i$ buffer, take a full decision tree
with $b$ measurements along each path. A node has its own fixed
position $u$ and branches on $T_{s,u}\in M$, with a separate subtree
or leaf for every value. Ancestors precede descendants; all nodes
can be ordered topologically. The resulting leaf history depends
only on the letters, not on the input vector.

After all these measurements, place aligned prediction trials with
fixed endpoints, checkpoints, and blocks, using the representation
$\rho$. Put
\[
              J'=|U_{i+1}|^2,\qquad
              R=|M|^{J'+1},\qquad b>R.
\]
For every ordered list of $R$ distinct leaf histories, allocate a trial
whose predictions across those histories exhaust $M^{J'+1}$. Give
arbitrary predictions to its other histories, and include a spare
trial with arbitrary predictions. Every history has a prediction on
every trial, including histories that cannot occur on an actual prefix.
Only predictions depend on history; physical trial intervals do not.

Define $g_{i+1,e}$ by the unshifted computation above, adding the ideal
shift of Lemma~\ref{lem:prediction} at each stage-$i$ trial end, using
the actual leaf history. These are its only shifts. Before $a_i$,
any buffers of other stages contribute only their letter actions
$\rho$; their shifts and terminal maps are not applied there. The
shifts are constants determined by the episode, so they do not change
its linear part. This is why the two splits can compute $g_{i+1}$
and then recover $g_i$ by Lemma~\ref{lem:affine-recovery} and the
basis-state calculation above.

\begin{lemma}[Total suffix tables]\label{lem:total-tables}
The tables $\Phi_i$ exist for $0\le i\le4$.
Every $g_{i+1,e}$ just defined is deterministic affine, independent of
split success and domain restrictions. Its linear part $L_{i+1,e}$
satisfies
\begin{equation}\label{eq:basis-recovery}
                 [v,1_M]L_{i+1,e}=[g_{i,e}(v),1_M].
\end{equation}
The state spaces, tree sizes, and aligned trial counts depend only on
the preceding state sizes and $M$, not on the length or complexity of
a table formula.
\end{lemma}

\begin{proof}
For $i=0$, all main boundaries and trials are after $a_0$. Their
marker decisions and actual interval products are suffix data, and
both clock graphs depend on fixed $q-t$. Replace each required
$T_{s,z}$ by $mT_{a_0,z}$. The explicit rules of Section~\ref{sec:outer}
then compute $F_1$, its homogeneous lift, and $F_2$ for every $m$,
agreeing with $F_2$ at the actual prefix product. This gives $\Phi_0$.

For $i<4$, put $a=a_{i+1}$ and $b=a_i$. Every stage-$i$ measurement
and trial lies after $a$. A hypothetical product $m$ assigns
$mT_{a,u}$ to the measurement at $u$, selecting a leaf $h(m)$ in
the full tree. Let $C_\nu$ be the stage-$i$ trial ends, and compute
their ideal shifts $\beta_\nu(h(m))$ from that leaf's predictions
and the actual interval products. For every $x\in U_{i+1}$, set
\begin{equation}\label{eq:total-table-formula}
 \Phi_{i+1}(m)(x)=
 \left(x\rho(mT_{a,b})+
       \sum_\nu\beta_\nu(h(m))\rho(T_{C_\nu,b})\right)H_i.
\end{equation}
The first term transports the input to $b$; each summand transports
one trial shift from its insertion point to $b$. Each shift is added
once to the whole vector, and the terminal map is then applied
linearly. The result is affine in $x$.

Every basis value encountered by $H_i$ has a defined image because
$\Phi_i$ is total on $M$. Hence this formula is defined even for
an infeasible hypothetical $m$ or an arbitrary sum of basis vectors.
At $m=T_{s,a}$, it is exactly $g_{i+1,e}$. A hypothetical $m$ changes
the work evolution and selected history; it changes neither the
actual suffix interval products nor $t,q$. This proves the next
table property.

For fixed word data, the letter actions and terminal map are linear
and all inserted shifts are independent of the input. Removing the
shifts therefore gives the linear part $\rho(T_{s,a_i})H_i$, proving
\eqref{eq:basis-recovery}. All schedules and basis sets are finite
functions of the already specified state sizes and $M$; formula
complexity changes none of them.
\end{proof}

The update and its total suffix table are now fixed independently of
successful splits. We next define the primary and secondary full sets:
they determine where each split must be available, and never redefine
the update that both splits compute.

\subsection{Copies and full sets}

There are two sorts of stencil kinds.
\begin{enumerate}[label=\textup{(\roman*)}]
\item A \emph{trial kind} specifies an aligned trial, its block, and a leaf
history. Its copies lie inside that block.
\item A \emph{measurement kind} specifies a tree node and an input payload
\(x\in U_{i+1}\). Its copies lie before that node's measurement and after all
its ancestor measurements.
\end{enumerate}
Every offset is distinct, so recovering an offset recovers all its metadata.

Inside \(D_i\), define the primary full set \(\mathcal F_i\) by
\begin{equation}\label{eq:primary-full}
 \text{the true history has a perfect aligned trial},\qquad
             \neg A_h^{\Gamma_i+\Delta_i}(t,q)\quad(h>i).
\end{equation}
Put \(D_i'=D_i\setminus\mathcal F_i\).
Inside \(D_i'\), define the secondary full set \(\mathcal O_i\) by
\begin{equation}\label{eq:secondary-full}
             \neg A_h^{\Gamma_i+\Delta_i+\Delta_i}(t,q)\quad(h>i),
\end{equation}
and put \(D_{i+1}=D_i'\setminus\mathcal O_i\).
Every episode left in \(D_{i+1}\) has a later witness expanded by
\(\Gamma_{i+1}\), proving the next entry promise.
At \(i=4\), the secondary condition is empty, so \(D_5=\varnothing\).

Both splits at this stage transmit the \emph{same} \(g_{i+1,e}\) from
Lemma~\ref{lem:total-tables}. These full sets and residual restrictions do
not redefine it.

\subsection{An exact suffix decoder}

At a cut \(k=s+p\), the suffix lists all possibilities
\begin{equation}\label{eq:time-candidates}
                \tau_{p'}=k+L_0-p',\qquad p'\in\mathcal P_i.
\end{equation}
In a primary test, it requires absence of every later
\(A_h^{\Gamma_i}\) at every candidate. In a secondary test, this guard is
strengthened to absence of every later \(A_h^{\Gamma_i+\Delta_i}\).

For \(i<4\), do the \(q\)-free last-type checks first and require equality of
\(\eta\) across all candidates. Apply the exact end rule at one reference
candidate, as in Lemma~\ref{lem:shared-anchor}, and use its \(q\).
Then perform the other required later-type checks. Finally require a unique
\(p'\in\mathcal P_i\) such that \(A_i^{\Gamma_i}(\tau_{p'},q)\).

At \(i=4\), there is no later guard. Use the uniqueness test for
\(A_4^{\Gamma_4}\), which is \(q\)-free, and only after decoding apply the
exact end rule at the decoded origin. Both uniqueness tests range over all
of \(\mathcal P_i\). After decoding, require metadata of the appropriate
kind. Every bounded control is mandatory; an argument missing any required
window rejects.

\begin{lemma}[Time decoder]\label{lem:time-decoder}
Paired with a genuine stencil prefix in the promised domain \(D_i\),
a successful primary or secondary decoder returns the actual offset and
consumes exactly the first episode.
On its respective full set, all later guards pass at every stencil cut,
and at most \(C_i\) offsets fail uniqueness. Thus every required kind has
a successfully decoded copy.
\end{lemma}

\begin{proof}
For an actual cut \(p\), the candidate \(p'=p\) is the true \(t\);
the others are \(t+p-p'\).

For \(i<4\), the explicit equality/end guard and \(h_i<w_E\) give the true
\(q\), independently of any decoded metadata. The true complete episode
contains every bounded control, and the proposed suffix must contain those
same controls. Thus even a prematurely or excessively long proposed suffix
sees the same control values.
The later guard eliminates every later witness at the true origin.
Since \(\Gamma_i\subseteq\Gamma_i+\Delta_i\), this holds for either guard.
The entry promise forces the true current witness to pass. Uniqueness
therefore selects precisely the actual \(p\), and the shared anchor already
gives exact consumption.

For \(i=4\), the entry promise itself makes the true candidate pass the
\(q\)-free test. Internal visibility of all controls prevents any choice of
proposed length or two-sided extension from hiding it. Uniqueness first
recovers the true time; applying its end rule then forces exact consumption.
This argument uses no anchor to prove the uniqueness assertion.

For availability of the primary guard, a candidate differs from \(t\) by
an element of \(\Delta_i\), and \eqref{eq:primary-full} excludes every
corresponding later \(\Gamma_i\)-translate.
For the secondary guard, \eqref{eq:secondary-full} excludes the further
\(\Delta_i\)-expanded tests at every candidate.
When \(i<4\), both full conditions include absence of \(A_4(t)\), since all
expansion sets contain zero. This is stability for \emph{every} shift of
magnitude at most \(w_E\), so the equality/end guard is available.

The true candidate now passes at each stencil cut. If actual offset \(p\)
is ambiguous, some \(p'\ne p\) satisfies the condition counted in
\eqref{eq:ambiguity-bound}. Choosing one such \(p'\) for each ambiguous \(p\)
injects ambiguous offsets into the counted pairs. At most \(C_i\) offsets
are lost in the whole stencil, so more than \(C_i\) copies of each kind
leave a decodable copy of that kind.
\end{proof}

For the secondary split one further consequence is crucial. On \(D_i'\),
a successful guard includes absence of every later
\(A_h^{\Gamma_i+\Delta_i}(t,q)\) at the true candidate.
If the true history had a perfect trial, the episode would satisfy
\eqref{eq:primary-full}, contradicting membership in \(D_i'\).
Thus the true history has no perfect trial. The same conclusion holds on
the secondary full set by \eqref{eq:secondary-full}.
The extra expansion is necessary here: membership in \(D_i'\) alone
allows an episode with a perfect true-history trial whose primary full
condition fails because of a later witness. Without the strengthened
guard, the secondary test could discard the actual measurement value.

The decoder has separated time recovery from state recovery. Any accepted
pair now knows the actual copy and its metadata. On each full set, the
whole-stencil count guarantees a copy of every required kind. We use
these facts first for a trial and then for a tree measurement.

\subsection{Primary transmission}

The primary \(P_x\) requires a trial-copy length \(p\). It checks its metadata
history against the actual measurements, processes from \(x\) to the current
trial start, and verifies the enumerated arrival vector, predicted transfer,
and preceding individual interval products of that block.
The whole tree and all earlier trials are on its side.

The primary \(Q_y\) applies the primary time decoder. Using the decoded
history, trial and block, it checks the suffix-product condition and computes
from the enumerated trial-end output through all later stage trials and the
terminal map, requiring result \(y\). Whole other trial shifts are evaluated
by ideal eligibility. In particular, the late map is evaluated on its stored
basis work values; the suffix does not need to learn a separate actual prefix
product at the cut.

For an accepted pair, Lemma~\ref{lem:time-decoder} identifies exactly the same
metadata on both sides. The trial conditions give ideal eligibility, and
Lemma~\ref{lem:prediction} proves \(g_{i+1,e}(x)=y\).
For a primary full episode, choose a perfect true-history trial and its
eligible block for the actually arriving vector. There is a decodable copy
of that history/trial/block kind, giving a split for every input vector.

\subsection{Secondary transmission and bad histories}

\begin{lemma}[Bad histories]\label{lem:bad-histories}
Fix one episode and the actual product tuples on all its aligned trial
intervals. Fewer than \(R\) leaf histories have no perfectly predicted trial,
including infeasible histories. Along the true path, fewer than \(R\) nodes
have a wrong proposed measurement value whose suffix continuation yields
such a history.
\end{lemma}

\begin{proof}
Let \(\mathcal B\) be the set of \emph{all} leaf histories whose prediction
differs from the actual tuple on every trial. If \(R\) distinct leaves
belonged to \(\mathcal B\), their allocated common trial would assign all
\(R\) possible tuples across those leaves. One assignment equals that trial's
actual tuple, a contradiction. Thus \(|\mathcal B|<R\).

At a true-path node with measurement \(u\), a wrong proposed value \(d\)
chooses a different branch. Continue at any later measurement \(u'\) using
\(dT_{u,u'}\). This selects a well-defined leaf in the off-path subtree,
whether or not the whole history is feasible.
Departures at two different nodes of the true path lie in disjoint subtrees,
so their leaves differ. Selecting one wrong surviving value at each ambiguous
node injects those nodes into \(\mathcal B\), proving the bound.
\end{proof}

The secondary \(P_x\) requires a measurement-copy length, verifies that its
node lies on the true path by checking ancestor measurements, and requires
payload \(x\).

The secondary \(Q_y\) applies the strengthened time decoder and obtains the
node, its ancestor-path metadata, and payload. At the node's measurement
position \(u\), try every \(d\in M\). Follow its branch and every later branch
using \(dT_{u,u'}\), and retain exactly the values whose resulting history
has no perfectly predicted trial. Require a unique retained value.

On an accepted pair in \(D_i'\), the true history has no perfect trial by
the strengthened-guard consequence above. Its actual value \(T_{s,u}\)
therefore survives, and uniqueness identifies it.
This datum suffices to compute the full update on every payload, not only on
a single basis vector. Indeed, if
\[
                        x=\sum_{v,m}c_{v,m}[v,m],
\]
then before \(u\) there have been no shifts in \(g_{i+1,e}\), and the vector
there is
\begin{equation}\label{eq:secondary-payload}
                x\rho(T_{s,u})
                   =\sum_{v,m}c_{v,m}[v,mT_{s,u}].
\end{equation}
The suffix knows all subsequent letters and trials, and all their ideal
shifts from the recovered history. It simulates to \(a_i\), applies the
total terminal map termwise, and requires result \(y\).
Each affine shift is added once to this entire vector, after the linear
work action; it is not added separately to each basis summand. The final
map \(H_i\) is applied once by linearity, and no letter action follows it.
No inverse in \(M\) is used.

For availability on a secondary full episode, the true history survives.
Lemma~\ref{lem:bad-histories} and \(b>R\) provide a true-path node at which no
wrong value survives. For any input payload, a decodable copy of its node
and payload kind exists. All continuation checks are after that cut, so this
copy gives an accepted split.

\subsection{The recursive conclusion}

\begin{lemma}[Aligned stage]\label{lem:aligned-stage}
Assume the entry promise, the table and placement properties, and
\eqref{eq:ambiguity-bound}, with more than \(C_i\) copies per kind.
Exact graph tests for \(g_{i+1}\) on \(D_{i+1}^*\) yield graph tests for
\(g_i\) on \(D_i^*\) by two uses of Lemma~\ref{lem:split} and finite Boolean
operations. All required hypotheses hold at the next residual stage, and
\(D_5=\varnothing\).
\end{lemma}

\begin{proof}
The secondary split first gives graph tests for \(g_{i+1}\) on \((D_i')^*\).
The primary split uses these as skips and gives its graph tests on \(D_i^*\).
Both splits act on \(U_{i+1}\) with the same deterministic affine update
\(g_{i+1}\). By
Lemma~\ref{lem:affine-recovery} these determine the product of its episode
linear parts. Equation~\eqref{eq:basis-recovery} shows that starting from
\([v,1_M]\), this product ends at \([g_{i,w}(v),1_M]\).
Thus finite basis-state queries recover the required graph tests for \(g_i\).

The domain promise and empty terminal domain follow from the definitions of
the full sets, and the next total table follows from
Lemma~\ref{lem:total-tables}. Starting with identity graph tests on the empty
\(D_5\) domain, the five stages consequently produce graph tests for
\(g_0=F_2\) on \(D_0^*\).
\end{proof}

